\documentclass[10pt,a4paper]{amsart}
\usepackage[margin=19mm]{geometry}
\usepackage{amsmath,amssymb,mathtools}
\usepackage[scr=rsfs,cal=euler]{mathalfa}
\usepackage{xfrac}
\usepackage[unicode,hidelinks,bookmarks=false]{hyperref}
\newcommand{\ie}{i.e.\ }
\newcommand{\cf}{cf.\ }
\newcommand{\resp}{respectively\ }
\newcommand{\Subsection}{\subsection}
\providecommand{\nobreakdash}{\nobreak\hbox{-}}
\setlength{\emergencystretch}{2em}
\multlinegap0pt
\usepackage{bm}
\usepackage{bbm}
\usepackage{dsfont}
\usepackage{xspace}
\usepackage{cancel}
\usepackage{mathtools}
\usepackage{amsthm}
\makeatletter
\@ifundefined{conjecture}{\newtheorem{conjecture}{Conjecture}}{}
\@ifundefined{lemma}{\newtheorem{lemma}[conjecture]{Lemma}}{}
\makeatother
\title[Rearrangements and infimum convolutions]{Rearrangements and infimum convolutions}
\author{Devraj Duggal, James Melbourne,, Cyril Roberto}
\date{Source: arXiv:2508.07983v1 (CC BY 4.0)}
\begin{document}
\maketitle

\noindent\emph{Source and licence.} Rearrangements and infimum convolutions. Version arXiv:2508.07983v1 (CC BY 4.0), distributed under the Creative Commons Attribution 4.0 International licence, as recorded in the arXiv licence metadata for this exact version and shown on the abstract page https://arxiv.org/abs/2508.07983v1. Full rights notice: licenses/arxiv-2508.07983-CC-BY-4.0.txt.\par\smallskip

\noindent\textit{Editorial scope.} Counted unit: the Lemma whose source label is \texttt{lem:absolutely-continuous} in the article (source0.tex lines 1064--1066, source label \texttt{lem:absolutely-continuous}), delivered together with the article's own complete original proof of it (source0.tex lines 1068--1137, one \texttt{proof} environment of 70 lines). No mathematics is written, repaired or re-derived: statement and proof are the article's own text line for line. Required context retained uncounted: none. Every definition, display and label of those lines is retained unchanged. Omitted from this package and not redistributed: 1--1063, 1067--1067, 1138--3664. Nothing inside a retained excerpt is omitted or altered: every retained body is delivered exactly as the article's lines are written, so no omission receipt of any kind accompanies this package. Citations: the retained excerpts carry no citation command at all, so no bibliography entry of the article is transcribed and no reference is redistributed. Reference adaptation: none was needed and none was made; every reference inside the delivered unit resolves inside it, so no unresolved reference or citation is produced. Printed slips kept literal: no printed word, symbol, hypothesis or number of the counted statement or of its proof was altered. Layout and class: the article's own class is \texttt{article} with packages including latexsym, color, xcolor, amsmath, amsthm, amssymb, amscd, amsfonts, stmaryrd, dsfont, graphicx, bbm, cancel, mathtools; the wrapper is the lane's pinned \texttt{amsart} editorial wrapper at 10pt on A4, which restates the theorem-like environments the retained text needs and adds the article's own macro definitions that the retained text needs (none), with extra wrapper packages (bm, bbm, dsfont, xspace, cancel, mathtools). The delivered numbering is the unit's own. Assets: no retained span contains a figure environment, a graphics-inclusion command, a PDF-page inclusion, a table or a TikZ picture; no image file is copied into this package, the package declares no asset dependency and no image file is redistributed with it. Licence: version arXiv:2508.07983v1 of this article is distributed under the Creative Commons Attribution 4.0 International licence, as recorded in the arXiv licence metadata for this exact version and shown on the abstract page https://arxiv.org/abs/2508.07983v1. A full rights notice is delivered at licenses/arxiv-2508.07983-CC-BY-4.0.txt. Characters: every retained excerpt is pure ASCII, so the delivered engine, XeLaTeX with its default Latin Modern text font, reports no missing character.
\begin{lemma}\label{lem:absolutely-continuous}
If $\mu_{n}\in\mathcal{S}$ converges weakly to a measure $\mu$,  then $\mu$ is absolutely continuous with respect to the Lebesgue measure. Moreover, if $e^{-\Psi_n}$ denotes the density of $\mu_n$ and $e^{-\Psi}$ the density of $\mu$, then $\Psi_n$ converges point-wise to $\Psi$ on the positive axis.
\end{lemma}

\begin{proof}
Denote by $e^{-\Psi_n}$ the density of $\mu_n$ with respect to the Lebesgue measure,
on $[0,\infty)$. Since $\mu_n \in \mathcal{S}$ we have $\mu_n(U) = \int_U e^{-\Psi_n(x)} dx \leq |U| e^{- \Psi_n(0)} = |U|$ for any Borel set $U$.  Thus for an open set $U$ such that $A \subseteq U$, by the weak convergence,
\begin{align*}
    \mu(A) \leq \mu(U) \leq \liminf_n \mu_n(U) \leq \liminf_n |U| = |U|.
\end{align*}
By the regularity of the Lebesgue measure, taking the infimum over all such $U$ gives,
\[
    \mu(A) \leq |A|.
\]
Thus $\mu$ has a density with respect to the Lebesgue measure, as expected. 

We will now show that $\lim_{n\rightarrow\infty}\Psi_{n}(x)=\Psi(x)$ for all $x \in [0,\infty)$.
By Helly's selection theorem we can extract from the sequence $(e^{-\Psi_n})$ a subsequence $(e^{-\Psi_{\varphi(n)}})$ converging point-wise to $e^{-\Psi}$.
In particular $\Psi$ is convex and therefore continuous on $(0,\infty)$.  
%Now every subsequence of $(\Psi_n(x))$ must converge to $\Psi(x)$, so that the whole sequence converge to $\Psi(x)$ and the proof is complete.



Remark that for all $a,b>0$, weak convergence implies that $\lim_{n\rightarrow\infty}\mu_{n}(a,b)=\mu(a,b)$. By contradiction, assume that $\limsup_{n\rightarrow\infty}e^{-\Psi_{n}(r)}>e^{-\Psi(r)}$ for some $r > 0$. 
The strict inequality guarantees  the existence of $\epsilon>0$ such that $\limsup_{n\rightarrow\infty}e^{-\Psi_{n}(r)} > \varepsilon + e^{-\Psi(r)}$. By monotonicity of $\Psi_n$, for any $0<a<r$ it holds
     \begin{align*}
      (r-a)(\varepsilon + e^{-\Psi(r)})
      < (r-a) \limsup_{n\rightarrow\infty} e^{-\Psi_n(r)} 
      \leq 
      \limsup_{n\rightarrow\infty}\mu_{n}(a,r)
      =\mu(a,r) 
      \leq (r-a)e^{-\Psi(a)} .
      \end{align*} 
      Therefore $\varepsilon + e^{-\Psi(r)} < e^{-\Psi(a)}$ that leads, by continuity of $\Psi$, in the limit $a \to r$, to a contradiction. Hence,  for all $r >0$, $\limsup_{n\rightarrow\infty}e^{-\Psi_{n}(r)} \leq e^{-\Psi(r)}$.
     
     Assume now that  $\liminf_{n\rightarrow\infty}e^{-\Psi_{n}(r)}<e^{-\Psi(r)}$ for some $r > 0$. Similarly, for some $\varepsilon>0$ and any $a>r$ it holds
  \begin{align*}
      (a-r)(-\varepsilon + e^{-\Psi(r)})
      > (a-r) \liminf_{n\rightarrow\infty} e^{-\Psi_n(r)} 
      \geq 
      \limsup_{n\rightarrow\infty}\mu_{n}(r,a)
      =\mu(r,a) 
      \geq (a-r)e^{-\Psi(a)} .
      \end{align*}    
This leads in the limit $a \to r$ to a contradiction and therefore to the desired conclusion: $\Psi_n(x) \to \Psi$ point-wise on $(0,\infty)$ and hence on $[0,\infty)$ since $\Psi_n(0)=\Psi(0)=0$.



% The strict inequality together with the continuity of $\Psi$ guarantee  the existence of $\epsilon,\delta>0$ such that $\limsup_{n\rightarrow\infty}\Psi_{n}(r)-\varepsilon>\Psi(r+\delta)$. By monotonicity of $\Psi_n$
%      \begin{align*}
%               \limsup_{n\rightarrow\infty}\mu_{n}(r,r+\delta)&\geq\delta\limsup_{n\rightarrow\infty}\Psi_{n}(r)\\
%              &>\delta\varepsilon+\delta\Psi(r+\delta)\\
%              &\geq\mu(r,r+\delta)
%       \end{align*} 
%       {\color{red}Cyril: the above computation is not correct.}
%      The case $\liminf_{n\rightarrow\infty}\Psi_{n}(r)<\Psi(r)$ may be treated analogously. 


%(but this density can be taken to be )bounded by $1$. 
%Curiously, we didn't use log-concavity like I suggested originally in our discussions!


% This is actually a general fact, if $\alpha$ is a fixed Radon measure, the operator $M_\alpha(\mu) = \left\| \frac{d\mu}{d\alpha} \right\|_\infty$ where the essential suprema is taken with respect to $\alpha$ and defined to be infinity when $\mu$ is not absolutely continuous with respect to $\alpha$, is lower semi-continuous by the same argument.  Even more generally, if  I recall correctly and the space of action is Polish, this can be taken as a statement about the lower semi-continuity of the $p$-R\'enyi divergence for $p > 0$ (this is the $p = \infty$ case above), or more generally still, any $f$-divergence being jointly lower semicontinuous (in both arguments simultaneously).

% Now denote $\frac{d\mu}{dx}=e^{-\Psi(x)}$. We will now show that $\lim_{n\rightarrow\infty}\Psi_{n}(x)=\Psi(x)$ for all $x \in [0,\infty)$.
% First remark that for $\forall a,b>0$, weak convergence implies that $\lim_{n\rightarrow\infty}\mu_{n}(a,b)=\mu(a,b)$. By contradiction, assume that $\limsup_{n\rightarrow\infty}\Psi_{n}(r)>\Psi(r)$ for some $r\geq 0$. The strict inequality guarantees  the existence of $\epsilon,\delta>0$ such that $\limsup_{n\rightarrow\infty}\Psi_{n}(r)-\varepsilon>\Psi(r+\delta)$. Observe the following
%     \begin{equation*}
%         \begin{split}
%             \limsup_{n\rightarrow\infty}\mu_{n}(r,r+\delta)&\geq\delta\limsup_{n\rightarrow\infty}\Psi_{n}(r)\\
%             &>\delta\varepsilon+\delta\Psi(r+\delta)\\
%             &\geq\mu(r,r+\delta)\\
%         \end{split}
%     \end{equation*} The case $\liminf_{n\rightarrow\infty}\Psi_{n}(r)<\Psi(r)$ may be treated analogously. 
\end{proof}

\end{document}
